\documentclass[10pt]{article}
\usepackage[margin=1in]{geometry}
\usepackage[T1]{fontenc}
\usepackage{lmodern,microtype,amsmath,amssymb,amsthm}
\theoremstyle{definition}
\newtheorem{assumption}{Assumption}
\begin{document}
\section*{Bounded Composite Saddle Optimization}
We consider $\min_{x\in X}\sup_{y\in Y}\Phi(x,y)$, where $\Phi(x,y)=g(x)+\langle y,\mathcal Kx\rangle-h(y)$ on real inner product spaces $E$ and $F$. The convex sets $X\subset E$ and $Y\subset F$ lie in radius-$R$ balls, with $R>0$. The functions $g:E\to\mathbb R$ and $h:F\to\mathbb R$ are convex and $M$-Lipschitz on their respective feasible sets, with $M\geq0$; the continuous linear map $\mathcal K:E\to F$ satisfies $\|\mathcal K\|\leq L$, with $L>0$. Each oracle call returns an exact operator action or constrained proximal point. Given $\varepsilon>0$, the output $(\bar x,\bar y)\in X\times Y$ must satisfy $\sup_{y\in Y}\Phi(\bar x,y)-\inf_{x\in X}\Phi(x,\bar y)\leq\varepsilon$ almost surely.

The spaces may be infinite dimensional. Their norms are induced by their inner products, and the product space is equipped with
\[
\|(x,y)\|_{E\times F}=\max\{\|x\|_E,\|y\|_F\}.
\]
The saddle problem is
\[
\min_{x\in X}\sup_{y\in Y}
\bigl\{g(x)+\langle y,\mathcal Kx\rangle-h(y)\bigr\}.
\]
The public inputs are $(E,F,X,Y,R,L,M,\varepsilon,x^0,y^0)$, where $x^0$ and $y^0$ are the initial points. The instance consists of $\mathcal K$, $g$, and $h$, accessed through the oracle described below.

\begin{assumption}[Feasible sets and parameters]
The sets $X$ and $Y$ are convex, and
\[
\begin{aligned}
x^0&\in X, & y^0&\in Y,\\
\|x\|_E&\leq R &&\text{for every }x\in X,\\
\|y\|_F&\leq R &&\text{for every }y\in Y,\\
R&>0, & L&>0,\qquad M\geq0,\qquad\varepsilon>0.
\end{aligned}
\]
Consequently, $X\times Y$ has diameter at most $2R$ in the product norm.
\end{assumption}

\begin{assumption}[Convex composite terms]
For every $t\in[0,1]$, $a,b\in X$, and $c,d\in Y$,
\[
\begin{aligned}
g((1-t)a+tb)&\leq(1-t)g(a)+tg(b),\\
h((1-t)c+td)&\leq(1-t)h(c)+th(d),\\
|g(a)-g(b)|&\leq M\|a-b\|_E,\\
|h(c)-h(d)|&\leq M\|c-d\|_F.
\end{aligned}
\]
\end{assumption}

\begin{assumption}[Linear coupling and proximal solutions]
The map $\mathcal K:E\to F$ is continuous and linear, with $\|\mathcal K\|\leq L$. There exists a map $\mathcal K^*:F\to E$ satisfying
\[
\langle\mathcal K^*y,z\rangle_E
=\langle y,\mathcal Kz\rangle_F
\qquad(y\in F,\ z\in E).
\]
There are maps $P_X:E\times\mathbb R\to E$ and $P_Y:F\times\mathbb R\to F$ such that, for every $\tau>0$, $u\in E$, and $v\in F$,
\[
\begin{aligned}
P_X(u,\tau)&\in X, & P_Y(v,\tau)&\in Y,\\
\tau g(P_X(u,\tau))
+\tfrac12\|P_X(u,\tau)-u\|_E^2
&\leq\tau g(w)+\tfrac12\|w-u\|_E^2
&& (w\in X),\\
\tau h(P_Y(v,\tau))
+\tfrac12\|P_Y(v,\tau)-v\|_F^2
&\leq\tau h(z)+\tfrac12\|z-v\|_F^2
&& (z\in Y).
\end{aligned}
\]
\end{assumption}

\begin{assumption}[Exact operator and proximal access]
Each call requests one of four operations: $\mathcal K^*y$ for any $y\in F$, $\mathcal Kx$ for any $x\in E$, $P_X(u,\tau)$ for any $(u,\tau)\in E\times\mathbb R$, or $P_Y(v,\tau)$ for any $(v,\tau)\in F\times\mathbb R$. The reply is the requested vector, exactly, at every call index and for every realization of oracle randomness. Each operation costs one oracle call. Proximal replies at $\tau\leq0$ are specified by these same maps but have no required minimization property. Queries may be adaptive and may lie outside the feasible sets. Any internal algorithm randomness is independent of oracle randomness.
\end{assumption}

The permitted algorithm is the Arrow--Hurwicz iteration with a fixed step size $\eta>0$, selected using only the public inputs. Starting from $(x^0,y^0)$, each complete iteration makes four calls in the order $\mathcal K^*y^t$, $P_X$, $\mathcal Kx^{t+1}$, and $P_Y$, implementing
\[
\begin{aligned}
x^{t+1}&=P_X(x^t-\eta\mathcal K^*y^t,\eta),\\
y^{t+1}&=P_Y(y^t+\eta\mathcal Kx^{t+1},\eta).
\end{aligned}
\]
Thus the coupling enters with the descent sign in $x$ and the ascent sign in $y$.

An externally specified budget $N\in\mathbb N$ may depend on the full public parameter tuple. Write $N=4n+r$, with $n\in\mathbb N$ and $r\in\{0,1,2,3\}$. The execution uses exactly $N$ oracle calls. If $n\geq1$, its output averages only the completed iterations:
\[
\bar x=\frac1n\sum_{t=1}^{n}x^t,
\qquad
\bar y=\frac1n\sum_{t=1}^{n}y^t.
\]
The remaining $r$ calls do not enter these averages. If $n=0$, the output is $(x^0,y^0)$ for $r=0,1$, and $(x^1,y^0)$ for $r=2,3$.

Writing $\mathbb P$ for the joint law of algorithm and oracle randomness, success requires
\[
(\bar x,\bar y)\in X\times Y,
\qquad
\sup_{y\in Y}\Phi(\bar x,y)
-\inf_{x\in X}\Phi(x,\bar y)
\leq\varepsilon
\qquad\mathbb P\text{-almost surely}.
\]
Equivalently, on a single event of probability one, the output is feasible and $\Phi(\bar x,y)-\Phi(x,\bar y)\leq\varepsilon$ holds simultaneously for every $x\in X$ and $y\in Y$.
\end{document}
